\chapter{Relative Spec}
\label{chap:Picard_RelativeSpec}
% archon:covers AlgebraicJacobian/Picard/RelativeSpec.lean

% STRATEGY NOTE (iter-172). This chapter is the foundational A.1.a sub-row of
% Route A (positive-genus arm of nonempty_jacobianWitness). It packages the
% relative-spectrum functor \underline{\Spec}_X(\mathcal{A}) : \mathrm{QcohAlg}(X)^{op}
% \to \mathrm{Sch}/X used by the relative Picard functor on a product C \times_k T.
% Mathlib has the absolute Spec functor and affine gluing (Scheme.GlueData,
% AffineScheme.glueOpens) but no general RelativeSpec construction; this chapter is
% the prover-ready specification for the new file AlgebraicJacobian/Picard/RelativeSpec.lean
% which iter-173 (file-skeleton lane) will scaffold.

\section{Setup and motivation}
\label{sec:relspec_setup}

The Route~A construction of \(J = \Pic^0_{C/k}\) as a \(g\)-dimensional abelian variety is
organised around the relative Picard functor
\(\Pic^\sharp_{C/k}(T) = \Pic(C \times_k T) / \pi_T^* \Pic(T)\)
on the category of \(k\)-schemes. Representing \(\Pic^\sharp_{C/k}\) requires the line-bundle
pullback construction on a relative curve \(C \times_k T \to T\), which in turn rests on a
\emph{relative scheme} construction on a quasi-coherent sheaf of \(\mathcal{O}_X\)-algebras
\(\mathcal{A}\), the \emph{relative spectrum}
\(\pi : \underline{\Spec}_X(\mathcal{A}) \to X\). This chapter packages the construction
of \(\underline{\Spec}_X(-)\) and its three core properties --- universal property, affine
base, base change --- as the foundational A.1.a building block. The relative Picard
functor itself (A.1.c), the line-bundle pullback functor on a product (A.1.b), and the
identity-component / dimension theory of \(J\) (A.3) live in separate sibling chapters
(\S~``Out of scope'' below).

% SOURCE: [Stacks Project], Section ``Relative spectrum via glueing'' and ``Relative
% spectrum as a functor'' (tags 01LL, 01LO, 01LQ, 01LR, 01LS, available locally at
% references/stacks-constructions.tex L309-L720). The Hartshorne reference is
% [Hartshorne], II Exercise 5.17 (read from references/hartshorne-algebraic-geometry.pdf).

\section{Quasi-coherent sheaves of algebras}
\label{sec:qc_sheaf_of_algebras}

\begin{definition}\leanok
[Quasi-coherent \(\mathcal{O}_X\)-algebra]
  \label{def:qc_sheaf_of_algebras}
  \lean{AlgebraicGeometry.Scheme.QcohAlgebra}
  % SOURCE: [Stacks Project], Section 01LL ``Relative spectrum via glueing'',
  % Situation ``situation-relative-spec''
  % (read from references/stacks-constructions.tex, L312-L318).
  % SOURCE QUOTE:
  % "Here \(S\) is a scheme, and \(\mathcal{A}\) is a quasi-coherent
  %  \(\mathcal{O}_S\)-algebra. This means that \(\mathcal{A}\) is a
  %  sheaf of \(\mathcal{O}_S\)-algebras which is quasi-coherent as an
  %  \(\mathcal{O}_S\)-module."
  \textit{Source: [Stacks Project], tag 01LL (situation-relative-spec).}
  Let \(X\) be a scheme. A \emph{quasi-coherent sheaf of \(\mathcal{O}_X\)-algebras} is a
  sheaf \(\mathcal{A}\) of \(\mathcal{O}_X\)-algebras (i.e. an \(\mathcal{O}_X\)-module
  equipped with \(\mathcal{O}_X\)-linear multiplication and unit morphisms satisfying
  associativity, commutativity, and the unit law as sheaf morphisms) whose underlying
  \(\mathcal{O}_X\)-module is quasi-coherent. We write \(\mathrm{QcohAlg}(X)\) for the
  category of quasi-coherent sheaves of \(\mathcal{O}_X\)-algebras and
  \(\mathcal{O}_X\)-algebra morphisms.
\end{definition}

% NOTE (iter-179 review): supersedes the iter-174 carrier note. The iter-179
% Block A refactor (analogies/relative-spec-encoding.md) upgrades the carrier
% to a 3-field structure
%     `structure QcohAlgebra (X : Scheme.{u}) where`
%     `  sheaf        : TopCat.Sheaf CommRingCat.{u} X.toPresheafedSpace`
%     `  unit         : X.sheaf ⟶ sheaf`
%     `  coequifibered : NatTrans.Coequifibered`
%     `    (Functor.whiskerLeft (AffineZariskiSite.toOpensFunctor X).op unit.hom)`
% The new third field IS the iter-175+ overlay the prior note flagged: Mathlib
% commit `b80f227`'s Stacks-01LL `NatTrans.Coequifibered` predicate is strictly
% weaker than the full `SheafOfModules.IsQuasicoherent` claimed by Definition
% 1.1's prose but is equivalent for the geometric purpose (constructing the
% relative-spectrum scheme) via the dense-subsite equivalence
% `AffineZariskiSite.sheafEquiv`. This is the predicate Mathlib's
% `AffineZariskiSite.relativeGluingData` consumes, so the iter-179
% `RelativeSpec`/`structureMorphism` bodies (now real Mathlib values) close
% against it.

\section{The relative-spectrum scheme}
\label{sec:relspec_scheme}

The construction is given affine-locally: on every affine open \(U = \Spec R \subseteq X\) one
sets \(\pi^{-1}(U) := \Spec(\mathcal{A}(U))\) regarded as an \(R\)-algebra, and the local pieces
glue compatibly because \(\mathcal{A}\) is quasi-coherent (the transition isomorphism
\(\mathcal{A}(U) \otimes_R S \xrightarrow{\sim} \mathcal{A}(V)\) for an inclusion
\(V = \Spec S \subseteq U\) gives an open immersion of the corresponding Specs).

\begin{theorem}\leanok
[Relative spectrum exists]
  \label{thm:relative_spec_exists}
  \lean{AlgebraicGeometry.Scheme.RelativeSpec}
  \uses{def:qc_sheaf_of_algebras}

  % SOURCE: [Stacks Project], tag 01LQ (lemma-glue-relative-spec)
  % (read from references/stacks-constructions.tex, L381-L405).
  % SOURCE QUOTE:
  % "In Situation REF.
  %  There exists a morphism of schemes
  %  \(\)
  %  \pi : \underline{\Spec}_S(\mathcal{A}) \longrightarrow S
  %  \(\)
  %  with the following properties:
  %  \begin{enumerate}
  %  \item for every affine open \(U \subset S\) there exists an isomorphism
  %  \(i_U : \pi^{-1}(U) \to \Spec(\mathcal{A}(U))\) over \(U\), and
  %  \item for \(U \subset U' \subset S\) affine open the composition
  %  \(\)
  %  \xymatrix{
  %  \Spec(\mathcal{A}(U)) \ar[r]^{i_U^{-1}} &
  %  \pi^{-1}(U) \ar[rr]^{inclusion} & &
  %  \pi^{-1}(U') \ar[r]^{i_{U'}} &
  %  \Spec(\mathcal{A}(U'))
  %  }
  %  \(\)
  %  is the open immersion of Lemma REF above.
  %  \end{enumerate}
  %  Moreover, \(\underline{\Spec}_S(\mathcal{A})\)
  %  is unique up to unique isomorphism over \(S\)."
  \textit{Source: [Stacks Project], tag 01LQ (lemma-glue-relative-spec);
  cf.\ [Hartshorne], II~Ex.~5.17(a).}
  Let \(X\) be a scheme and \(\mathcal{A}\) a quasi-coherent \(\mathcal{O}_X\)-algebra
  (\cref{def:qc_sheaf_of_algebras}). There exists a scheme
  \(\underline{\Spec}_X(\mathcal{A})\) together with an affine morphism
  \(\pi : \underline{\Spec}_X(\mathcal{A}) \to X\) such that
  \begin{enumerate}
    \item for every affine open \(U \subseteq X\), there is an isomorphism
      \(i_U : \pi^{-1}(U) \xrightarrow{\sim} \Spec(\mathcal{A}(U))\) over \(U\);
    \item for \(U \subseteq U' \subseteq X\) both affine open, the composite
      $\Spec(\mathcal{A}(U)) \xrightarrow{i_U^{-1}} \pi^{-1}(U) \hookrightarrow
      \pi^{-1}(U') \xrightarrow{i_{U'}} \Spec(\mathcal{A}(U'))$
      coincides with the open immersion induced by the restriction
      \(\mathcal{A}(U') \to \mathcal{A}(U)\).
  \end{enumerate}
  The pair \((\underline{\Spec}_X(\mathcal{A}), \pi)\) is unique up to unique isomorphism
  over \(X\).
\end{theorem}

% SOURCE QUOTE PROOF: "Follows immediately from
% Lemmas REF,
% REF, and
% REF.
% Uniqueness is stated in the last sentence of
% Lemma REF."
\begin{proof}
  
  
\leanok
  By the open-cover gluing principle for schemes (Mathlib
  \texttt{Scheme.GlueData} / \texttt{AffineScheme.glueOpens}) it suffices to specify
  affine local pieces \(\Spec(\mathcal{A}(U))\) for each affine open \(U \subseteq X\) and
  open-immersion transition data
  \(\Spec(\mathcal{A}(U)) \hookrightarrow \Spec(\mathcal{A}(U'))\) for \(U \subseteq U'\)
  satisfying the cocycle condition. The transition morphisms come from the ring maps
  \(\mathcal{A}(U') \to \mathcal{A}(U)\); quasi-coherence yields the cartesian square
  \(\Spec(\mathcal{A}(U)) = U \times_{U'} \Spec(\mathcal{A}(U'))\), so each transition is an
  open immersion. Transitivity (Stacks lemma-transitive-spec) gives the cocycle condition.
  Uniqueness of the glued scheme is the standard property of the gluing functor.
\end{proof}

\begin{definition}

\leanok
[Structure morphism of the relative spectrum]
  \label{def:relspec_structure_morphism}
  \lean{AlgebraicGeometry.Scheme.RelativeSpec.structureMorphism}
  \uses{thm:relative_spec_exists, def:qc_sheaf_of_algebras}
  Let \(X\) be a scheme and \(\mathcal{A}\) a quasi-coherent \(\mathcal{O}_X\)-algebra.
  The \emph{structure morphism}
  \(\pi : \underline{\Spec}_X(\mathcal{A}) \to X\) is the canonical projection of the
  relative spectrum to its base, obtained as the colimit descent of the affine-local
  structure morphisms \(\Spec(\mathcal{A}(U)) \to U\) over the directed affine open
  cover of \(X\). It is the morphism whose affineness and base-change behaviour are the
  subject of the universal-property and base-change results below.
\end{definition}

\begin{proof}

\leanok
  Constructed directly in Lean as the base morphism of the relative gluing data of
  \(\mathcal{A}\); this is the projection underlying \cref{thm:relative_spec_exists}.
\end{proof}

\begin{theorem}[Universal property of the relative spectrum]
  \label{thm:relative_spec_univ}
  \lean{AlgebraicGeometry.Scheme.RelativeSpec.UniversalProperty}
  % NOTE iter-283: removed the spurious `thm:relative_spec_base_change` from this
  % \uses. The universal property is the PRIMITIVE statement, proved by reducing
  % to the affine-base case and gluing (Lean body of `UniversalProperty` uses
  % only `isAffineHom_of_forall_exists_isAffineOpen`/`relativeGluingData`, never
  % `base_change`). Base change is a DOWNSTREAM corollary
  % (thm:relative_spec_base_change \uses thm:relative_spec_univ), so listing it
  % here created a univ <-> base_change dependency cycle that crashed
  % `leanblueprint web`.
  \uses{thm:relative_spec_exists, thm:relative_spec_affine_base,
    def:relspec_structure_morphism}

  % NOTE (iter-173 review): the iter-173 file-skeleton encodes this as the
  % structural consequence ``IsAffineHom (structureMorphism 𝒜)'' rather than
  % the natural-bijection statement pinned by this prose. Iter-174+ must
  % upgrade the Lean type to a ``CategoryTheory.Functor.RepresentableBy''
  % witness (or equivalent) before any consumer file relies on this theorem.
  % SOURCE: [Stacks Project], tag 01LQ continuation: Section ``Relative spectrum as
  % a functor'', functor F and Lemma lemma-spec
  % (read from references/stacks-constructions.tex, L427-L466 + L547-L551).
  % SOURCE QUOTE:
  % "For any \(f : T \to S\) the pullback
  %  \(f^*\mathcal{A}\) is a quasi-coherent sheaf of \(\mathcal{O}_T\)-algebras.
  %  We are going to consider pairs \((f : T \to S, \varphi)\) where
  %  \(f\) is a morphism of schemes and \(\varphi : f^*\mathcal{A} \to \mathcal{O}_T\)
  %  is a morphism of \(\mathcal{O}_T\)-algebras. Note that this is the
  %  same as giving a \(f^{-1}\mathcal{O}_S\)-algebra homomorphism
  %  \(\varphi : f^{-1}\mathcal{A} \to \mathcal{O}_T\), see
  %  Sheaves, Lemma REF.
  %  This is also the same as giving an \(\mathcal{O}_S\)-algebra map
  %  \(\varphi : \mathcal{A} \to f_*\mathcal{O}_T\), see
  %  Sheaves, Lemma REF.
  %  We will use all three ways of thinking about \(\varphi\),
  %  without further mention. [...]
  %  In this way we have defined a functor
  %  \begin{eqnarray}
  %  \label{equation-spec}
  %  F : \Sch^{opp} & \longrightarrow & \textit{Sets} \\
  %  T & \longmapsto & F(T) = \{\text{pairs }(f, \varphi) \text{ as above}\}
  %  \nonumber
  %  \end{eqnarray}
  %  [...]
  %  The functor \(F\) is representable by a scheme."
  \textit{Source: [Stacks Project], tag 01LQ (lemma-spec + section-spec
  functor description); [Hartshorne], II~Ex.~5.17(b).}
  Let \(X\) be a scheme and \(\mathcal{A}\) a quasi-coherent \(\mathcal{O}_X\)-algebra.
  Then \(\underline{\Spec}_X(\mathcal{A})\) represents the functor
  \(F : \Sch^{op} \to \mathrm{Set}\),
  $T \mapsto F(T) := \{(f, \varphi) : f : T \to X \text{ a morphism},\;
  \varphi : f^* \mathcal{A} \to \mathcal{O}_T \text{ an } \mathcal{O}_T\text{-algebra map}\}.$
  Equivalently, for any \(X\)-scheme \(g : T \to X\),
  there is a natural bijection
  \[
    \Hom_X(T, \underline{\Spec}_X(\mathcal{A})) \;\;\cong\;\;
    \Hom_{\mathcal{O}_X\text{-alg}}(\mathcal{A}, g_* \mathcal{O}_T)
  \]
  (using the adjunction between \(g_*\) and \(g^*\) to convert between
  \(\mathcal{A} \to g_*\mathcal{O}_T\) and \(g^*\mathcal{A} \to \mathcal{O}_T\)).
  The bijection is natural in \(T\).
\end{theorem}

% SOURCE QUOTE PROOF: TODO retrieve from references/stacks-constructions.tex
% (proof of lemma-spec, L553-L600, ~50 lines covering the Zariski sheaf property,
% the subfunctors F_i associated to an affine cover, and identifying each F_i
% with the affine-base case via lemma-spec-base-change and lemma-spec-affine).
% Verbatim quote omitted for length; the structural steps are restated below in
% project notation.
\begin{proof}
  
  % NOTE (iter-179 review): supersedes the iter-176 placeholder-discharge note.
  % The iter-179 Block A refactor replaces the iter-176 placeholder bodies with
  % the Mathlib values ``(AffineZariskiSite.relativeGluingData _).glued`` /
  % ``.toBase``; iter-179 Lane B closes the Lean ``UniversalProperty`` proof
  % via ``isAffineHom_of_forall_exists_isAffineOpen`` +
  % ``relativeGluingData.toBase_preimage_eq_opensRange_ι`` +
  % ``isAffineOpen_opensRange`` — a real proof against the gluing infrastructure,
  % NOT against a trivial identity placeholder. The Lean type remains the
  % strictly weaker ``IsAffineHom (structureMorphism 𝒜)'' rather than the
  % Yoneda-bijection form pinned above (iter-174+ commitment; still pending
  % iter-180+). Consumers may now rely on the affineness conclusion, but the
  % Yoneda-bijection witness is still deferred.
  Following Stacks tag 01LQ (proof of lemma-spec): the functor \(F\) is a Zariski sheaf
  because morphisms of schemes and \(\mathcal{O}\)-algebra maps both glue under open covers.
  Choose an affine open cover \(X = \bigcup U_i\). The subfunctor
  \(F_i \subseteq F\) of pairs \((f, \varphi)\) with \(f(T) \subseteq U_i\) is representable by
  \(\Spec(\mathcal{A}(U_i))\) via the affine-base case
  (\cref{thm:relative_spec_affine_base}) together with the
  base-change identification \(F_i \cong F_{U_i,\,\mathcal{A}|_{U_i}}\)
  (\cref{thm:relative_spec_base_change}). Each \(F_i \hookrightarrow F\) is an open
  subfunctor (pull back \(U_i\) along \(f\)). The \(F_i\) jointly cover \(F\) since the \(U_i\)
  cover \(X\). The gluing of representable open subfunctors produces a representing scheme,
  which is canonically isomorphic to the glued scheme \(\underline{\Spec}_X(\mathcal{A})\)
  of \cref{thm:relative_spec_exists}.
\end{proof}

\section{Affine base}
\label{sec:relspec_affine}

\begin{theorem}[Relative spectrum over an affine base]
  \label{thm:relative_spec_affine_base}
  \lean{AlgebraicGeometry.Scheme.RelativeSpec.affine_base_iff}
  \uses{def:qc_sheaf_of_algebras}
  % NOTE (iter-173 review): the iter-173 file-skeleton encodes this as the
  % weaker ``IsAffine ((Spec R).RelativeSpec 𝒜)'' rather than the canonical
  % iso ``Spec_X(𝒜) ≅ Spec(Γ(X, 𝒜))'' pinned by this prose. The Lean name
  % ``affine_base_iff'' is also misleading (the encoded type is not an iff).
  % Iter-174+ refines to the iff-with-Γ statement once Γ : QcohAlgebra → CommRingCat
  % is available.
  
  % SOURCE: [Stacks Project], tag 01LO (lemma-spec-affine)
  % (read from references/stacks-constructions.tex, L491-L545).
  % SOURCE QUOTE:
  % "In Situation REF.
  %  Let \(F\) be the functor associated to \((S, \mathcal{A})\) above.
  %  If \(S\) is affine, then \(F\) is representable by the
  %  affine scheme \(\Spec(\Gamma(S, \mathcal{A}))\)."
  \textit{Source: [Stacks Project], tag 01LO (lemma-spec-affine).}
  Let \(X = \Spec R\) be an affine scheme and let \(\mathcal{A}\) be a quasi-coherent
  \(\mathcal{O}_X\)-algebra. Set \(A := \Gamma(X, \mathcal{A})\), regarded as an
  \(R\)-algebra. Then there is a canonical isomorphism of \(X\)-schemes
  \[
    \underline{\Spec}_X(\mathcal{A}) \;\;\cong\;\; \Spec(A).
  \]
  In particular, on an affine base the relative spectrum reduces to the absolute
  spectrum of the global sections.
\end{theorem}

% SOURCE QUOTE PROOF: "Write \(S = \Spec(R)\) and \(A = \Gamma(S, \mathcal{A})\).
% Then \(A\) is an \(R\)-algebra and \(\mathcal{A} = \widetilde A\).
% The ring map \(R \to A\) gives rise to a canonical map
% \(\)
% f_{univ} : \Spec(A)
% \longrightarrow
% S = \Spec(R).
% \(\)
% We have
% \(f_{univ}^*\mathcal{A} =  \widetilde{A \otimes_R A}\)
% by Schemes, Lemma REF.
% Hence there is a canonical map
% \(\)
% \varphi_{univ} :
% f_{univ}^*\mathcal{A} = \widetilde{A \otimes_R A}
% \longrightarrow
% \widetilde A = \mathcal{O}_{\Spec(A)}
% \(\)
% coming from the \(A\)-module map \(A \otimes_R A \to A\),
% \(a \otimes a' \mapsto aa'\). We claim that the pair
% \((f_{univ}, \varphi_{univ})\) represents \(F\) in this case."
\begin{proof}
  
  % NOTE (iter-179 review): supersedes the iter-176 placeholder-discharge note.
  % The iter-176 ``change IsAffine (Spec R); inferInstance`` form is now stale
  % — iter-179 Lane B closes via ``UniversalProperty 𝒜`` + ``isAffine_of_isAffineHom``
  % against the iter-179 Mathlib-backed ``RelativeSpec``/``structureMorphism``
  % bodies. The Lean type remains the strictly weaker ``IsAffine`` rather than
  % the canonical iso ``Spec_X(𝒜) ≅ Spec(Γ(X,𝒜))'' named below (iter-174+
  % commitment; still pending iter-180+).
  Write \(X = \Spec R\) and \(A = \Gamma(X, \mathcal{A})\). Quasi-coherence of
  \(\mathcal{A}\) gives \(\mathcal{A} = \widetilde{A}\) as \(\mathcal{O}_X\)-modules
  (with the \(R\)-algebra structure on \(A\) inducing the algebra structure on
  \(\widetilde A\)). The ring map \(R \to A\) gives a morphism
  \(f_{\mathrm{univ}} : \Spec A \to \Spec R = X\) and a canonical
  \(\mathcal{O}_{\Spec A}\)-algebra map
  $\varphi_{\mathrm{univ}} : f_{\mathrm{univ}}^* \mathcal{A} = \widetilde{A \otimes_R A}
  \to \mathcal{O}_{\Spec A} = \widetilde A$ given by multiplication
  \(a \otimes a' \mapsto aa'\). We verify \((\Spec A, f_{\mathrm{univ}}, \varphi_{\mathrm{univ}})\)
  represents \(F\) on the affine base: given a test scheme \(T\) and a pair
  \((f, \varphi)\), the composite
  $A = \Gamma(X, \mathcal{A}) \xrightarrow{f^\sharp}
  \Gamma(T, f^*\mathcal{A}) \xrightarrow{\varphi} \Gamma(T, \mathcal{O}_T)$
  is a ring map \(A \to \Gamma(T, \mathcal{O}_T)\) and hence corresponds to a morphism
  \(T \to \Spec A\) by the affine universal property
  (Mathlib \texttt{Scheme.toSpec\_naturality} /
  \texttt{Scheme.Hom.toSpec}). One checks this assignment inverts the universal map
  \((f, \varphi) \mapsto a\) defined by \(f = f_{\mathrm{univ}} \circ a\),
  \(\varphi = a^* \varphi_{\mathrm{univ}}\).
\end{proof}

\section{Base change}
\label{sec:relspec_basechange}

\begin{theorem}[Base change of the relative spectrum]
  \label{thm:relative_spec_base_change}
  \lean{AlgebraicGeometry.Scheme.RelativeSpec.base_change}
  \uses{thm:relative_spec_univ, lem:relspec_pullback_iso,
    lem:relspec_pullback_coequifibered}

  % NOTE (iter-173 review): the iter-173 file-skeleton encodes this as the
  % weaker existential ``∃ 𝒜', Nonempty (pullback ≅ T.RelativeSpec 𝒜')'',
  % rather than the canonical iso with a *named* pullback ``g^* 𝒜'' pinned
  % by this prose. Iter-174+ factors out ``pullbackQcoh g 𝒜'' and tightens
  % the statement to the canonical iso.
  % SOURCE: [Stacks Project], tag 01LS (lemma-spec-base-change and
  % lemma-spec-properties part (2))
  % (read from references/stacks-constructions.tex, L467-L489 + L662-L691).
  % SOURCE QUOTE:
  % "In Situation REF.
  %  Let \(F\) be the functor
  %  associated to \((S, \mathcal{A})\) above.
  %  Let \(g : S' \to S\) be a morphism of schemes.
  %  Set \(\mathcal{A}' = g^*\mathcal{A}\). Let \(F'\) be the
  %  functor associated to \((S', \mathcal{A}')\) above.
  %  Then there is a canonical isomorphism
  %  \(\)
  %  F' \cong h_{S'} \times_{h_S} F
  %  \(\)
  %  of functors. [...]
  %  For every morphism \(g : S' \to S\) we have
  %  $S' \times_S \underline{\Spec}_S(\mathcal{A}) =
  %  \underline{\Spec}_{S'}(g^*\mathcal{A})$."
  \textit{Source: [Stacks Project], tag 01LS (lemma-spec-base-change +
  lemma-spec-properties).}
  Let \(g : T \to X\) be a morphism of schemes and let \(\mathcal{A}\) be a quasi-coherent
  \(\mathcal{O}_X\)-algebra. Then \(g^* \mathcal{A}\) is a quasi-coherent
  \(\mathcal{O}_T\)-algebra and there is a canonical isomorphism of \(T\)-schemes
  \[
    T \times_X \underline{\Spec}_X(\mathcal{A}) \;\;\cong\;\;
    \underline{\Spec}_T(g^* \mathcal{A}).
  \]
\end{theorem}

% SOURCE QUOTE PROOF: "A pair \((f' : T \to S', \varphi' : (f')^*\mathcal{A}' \to \mathcal{O}_T)\)
% is the same as a pair \((f, \varphi : f^*\mathcal{A} \to \mathcal{O}_T)\)
% together with a factorization of \(f\) as \(f = g \circ f'\). Namely with
% this notation we have
% \((f')^* \mathcal{A}' = (f')^*g^*\mathcal{A} = f^*\mathcal{A}\).
% Hence the lemma."
\begin{proof}
  
  % NOTE (iter-185 review, supersedes iter-179): the iter-185 Lane D HARD-BAR
  % closed BOTH iter-179 Mathlib-gap sorries kernel-clean:
  % (i) ``Scheme.QcohAlgebra.pullback_coequifibered`` (the ``coequifibered``
  % field, RelativeSpec.lean L358-377) via
  % ``coequifibered_iff_forall_isLocalizationAway`` +
  % ``IsAffineOpen.isLocalization_of_eq_basicOpen``;
  % (ii) ``Scheme.RelativeSpec.pullback_iso_desc_isIso`` (the canonical iso body,
  % RelativeSpec.lean L546-632) via the explicit 3-iso chain
  % (hPre identity \(\to\) isoOpensRange.symm \(\to\) pullback_iso_affine_piece.symm)
  % and ``IsZariskiLocalAtTarget`` affineness criterion.
  % The Lean type remains the weaker existential rather than the canonical iso
  % with a *named* pullback ``g^*𝒜'' pinned by this prose (iter-174+ signature
  % refinement still pending — see chapter NOTE on the statement block above).
  % Consumers may now consume ``pullback_iso''s witness as a fully axiom-clean
  % closure (no more "load-bearing-pending" qualifier); A.1.a body-level work
  % is functionally complete.
  By \cref{thm:relative_spec_univ}, \(\underline{\Spec}_X(\mathcal{A})\) represents
  the functor \(F\) sending an \(X\)-scheme \(h : S \to X\) to the set of
  \(\mathcal{O}_S\)-algebra maps \(h^*\mathcal{A} \to \mathcal{O}_S\). A morphism
  \(S \to T \times_X \underline{\Spec}_X(\mathcal{A})\) over \(T\) is the same as a pair of an
  \(X\)-morphism \(S \to \underline{\Spec}_X(\mathcal{A})\) together with a \(T\)-factorisation
  of \(S \to X\), i.e.\ a morphism \(S \to T\) whose composition with \(g\) equals
  \(S \to X\). Tracing through, this is a \(T\)-morphism \(f' : S \to T\) together with an
  \(\mathcal{O}_S\)-algebra map \((f')^*(g^*\mathcal{A}) = f^*\mathcal{A} \to \mathcal{O}_S\).
  This is precisely \(F'(S)\), the functor for \((T, g^*\mathcal{A})\), so the canonical map
  \(T \times_X \underline{\Spec}_X(\mathcal{A}) \to \underline{\Spec}_T(g^*\mathcal{A})\)
  is an isomorphism by Yoneda.
\end{proof}

\section{Pullback compatibility helpers}
\label{sec:relspec_pullback_helpers}

This section records the substrate lemmas and constructions underlying the
base-change isomorphism (\cref{thm:relative_spec_base_change}). Throughout,
\(g : T \to X\) is a morphism of schemes, \(\mathcal{A}\) is a quasi-coherent
\(\mathcal{O}_X\)-algebra, and \(q\) denotes the first projection of the pullback
\(T \times_X \underline{\Spec}_X(\mathcal{A})\) along the structure morphism
\(\pi : \underline{\Spec}_X(\mathcal{A}) \to X\). The pulled-back quasi-coherent
\(\mathcal{O}_T\)-algebra \(g^*\mathcal{A}\) is realised as the pushforward
\(q_* \mathcal{O}_{T \times_X \underline{\Spec}_X(\mathcal{A})}\).

\begin{lemma}
[Affineness of the structural pullback projection]
  \label{lem:relspec_pullback_fst_isaffinehom}
  \lean{AlgebraicGeometry.Scheme.QcohAlgebra.pullback_fst_isAffineHom}
  \uses{thm:relative_spec_univ, def:relspec_structure_morphism}
  The first projection
  \(q : T \times_X \underline{\Spec}_X(\mathcal{A}) \to T\) of the pullback of the
  structure morphism \(\pi\) along \(g\) is an affine morphism.
\end{lemma}

\begin{proof}
  The structure morphism \(\pi\) is affine by \cref{thm:relative_spec_univ}, and
  affineness of morphisms is stable under base change; hence its pullback \(q\) is
  affine. Proved directly in Lean.
\end{proof}

\begin{lemma}
[Quasi-coherence overlay for the pushforward along the affine projection]
  \label{lem:relspec_pullback_coequifibered}
  \lean{AlgebraicGeometry.Scheme.QcohAlgebra.pullback_coequifibered}
  \uses{lem:relspec_pullback_fst_isaffinehom}
  The pushforward \(q_* \mathcal{O}_{T \times_X \underline{\Spec}_X(\mathcal{A})}\)
  along the affine projection \(q\) carries the Stacks~01LL quasi-coherence overlay:
  for every affine open \(U \subseteq T\) and section \(r \in \Gamma(T, U)\), the
  restriction of the pushforward to the basic open \(D(r) \subseteq U\) is the
  localization away from \(r\) of its sections over \(U\). This is the predicate making
  \(g^*\mathcal{A}\) a well-formed quasi-coherent \(\mathcal{O}_T\)-algebra.
\end{lemma}

\begin{proof}
  Since \(q\) is affine (\cref{lem:relspec_pullback_fst_isaffinehom}), the preimage
  \(q^{-1}(U)\) of an affine open \(U\) is affine, and \(q^{-1}(D(r))\) is the basic open
  of \(q^{-1}(U)\) cut out by the pulled-back section; the localization-away
  identification then follows from the affine-open localization criterion. Proved
  directly in Lean.
\end{proof}

\begin{definition}

[Base change of a quasi-coherent algebra]
  \label{def:relspec_qcoh_pullback}
  \lean{AlgebraicGeometry.Scheme.QcohAlgebra.pullback}
  \uses{def:qc_sheaf_of_algebras, lem:relspec_pullback_fst_isaffinehom,
    lem:relspec_pullback_coequifibered}
  Let \(g : T \to X\) be a morphism of schemes and \(\mathcal{A}\) a quasi-coherent
  \(\mathcal{O}_X\)-algebra (\cref{def:qc_sheaf_of_algebras}). The \emph{base change}
  (pullback) \(g^*\mathcal{A}\) of \(\mathcal{A}\) along \(g\) is the quasi-coherent
  \(\mathcal{O}_T\)-algebra realised as the pushforward
  \(q_* \mathcal{O}_{T \times_X \underline{\Spec}_X(\mathcal{A})}\) along the first
  projection \(q : T \times_X \underline{\Spec}_X(\mathcal{A}) \to T\) of the pullback of
  the structure morphism \(\pi\) along \(g\). Concretely, the underlying
  \(\mathcal{O}_T\)-module is the topological pushforward \(q_* \mathcal{O}\) and the unit
  is the comorphism \(q^\sharp : \mathcal{O}_T \to q_* \mathcal{O}\); the Stacks~01LL
  quasi-coherence (coequifibered) overlay is supplied by
  \cref{lem:relspec_pullback_coequifibered}, which itself rests on the affineness of \(q\)
  (\cref{lem:relspec_pullback_fst_isaffinehom}). This is the project's base-change
  constructor for quasi-coherent algebras, packaging \(g^*\mathcal{A}\) as a value of
  \(\mathrm{QcohAlg}(T)\); it is the object underlying the base-change isomorphism
  (\cref{thm:relative_spec_base_change}).
\end{definition}

\begin{proof}

  Proved directly in Lean.
\end{proof}

\begin{definition}
[Per-affine-open piece of the base-change iso]
  \label{def:relspec_pullback_iso_affine_piece}
  \lean{AlgebraicGeometry.Scheme.RelativeSpec.pullback_iso_affine_piece}
  \uses{lem:relspec_pullback_fst_isaffinehom}
  For an affine open \(U \subseteq T\), the preimage \(q^{-1}(U)\) is affine and is
  canonically isomorphic to \(\Spec\) of the sections of the pulled-back algebra
  \(g^*\mathcal{A}\) over \(U\). This per-piece isomorphism feeds the global base-change
  iso construction.
\end{definition}

\begin{proof}
  As \(q\) is affine (\cref{lem:relspec_pullback_fst_isaffinehom}), \(q^{-1}(U)\) is
  affine; its canonical isomorphism with \(\Spec \Gamma(q^{-1}(U), \mathcal{O})\) matches
  the sections of the pushforward \(g^*\mathcal{A}\) over \(U\) by definitional
  unfolding. Constructed directly in Lean.
\end{proof}

\begin{definition}
[Pullback cocone on the relative-gluing functor]
  \label{def:relspec_pullback_cocone}
  \lean{AlgebraicGeometry.Scheme.RelativeSpec.pullback_cocone}
  \uses{lem:relspec_pullback_fst_isaffinehom}
  A cocone on the relative-gluing-data functor of \(g^*\mathcal{A}\) whose apex is the
  pullback \(T \times_X \underline{\Spec}_X(\mathcal{A})\); its component at an affine
  open \(U\) is the canonical map \(\Spec((g^*\mathcal{A})(U)) \to
  T \times_X \underline{\Spec}_X(\mathcal{A})\) into the affine piece \(q^{-1}(U)\).
\end{definition}

\begin{proof}
  The components are the affine-open inclusion maps of the pulled-back pieces; their
  compatibility (naturality) follows from the functoriality of these inclusions under
  restriction along basic opens. Constructed directly in Lean.
\end{proof}

\begin{lemma}
[Descent of the pullback cocone covers the structure morphism]
  \label{lem:relspec_pullback_cocone_desc_comp_fst}
  \lean{AlgebraicGeometry.Scheme.RelativeSpec.pullback_cocone_desc_comp_fst}
  \uses{def:relspec_pullback_cocone, lem:relspec_pullback_fst_isaffinehom}
  The colimit descent of the pullback cocone
  (\cref{def:relspec_pullback_cocone}), as a morphism
  \(\underline{\Spec}_T(g^*\mathcal{A}) \to T \times_X \underline{\Spec}_X(\mathcal{A})\),
  composed with the projection \(q\), equals the structure morphism
  \(\underline{\Spec}_T(g^*\mathcal{A}) \to T\) of the relative spectrum of
  \(g^*\mathcal{A}\).
\end{lemma}

\begin{proof}
  By the universal property of the colimit it suffices to check the identity on each
  cocone component, where it reduces to the compatibility of the affine-piece
  inclusion with the projection \(q\). Proved directly in Lean.
\end{proof}

\begin{lemma}
[The base-change descent is an isomorphism]
  \label{lem:relspec_pullback_iso_desc_isiso}
  \lean{AlgebraicGeometry.Scheme.RelativeSpec.pullback_iso_desc_isIso}
  \uses{def:relspec_pullback_cocone, lem:relspec_pullback_cocone_desc_comp_fst,
    def:relspec_pullback_iso_affine_piece, lem:relspec_pullback_fst_isaffinehom}
  The colimit descent map
  \(\underline{\Spec}_T(g^*\mathcal{A}) \to T \times_X \underline{\Spec}_X(\mathcal{A})\)
  determined by the pullback cocone is an isomorphism of schemes.
\end{lemma}

\begin{proof}
  Being an isomorphism is local on the target, so it suffices to check it over each
  member \(q^{-1}(U)\) of the affine open cover. There the descent restricts, via
  \cref{lem:relspec_pullback_cocone_desc_comp_fst} and the range identification of the
  gluing inclusions, to the per-affine-open isomorphism
  (\cref{def:relspec_pullback_iso_affine_piece}). Proved directly in Lean.
\end{proof}

\begin{definition}
[Base-change iso construction]
  \label{def:relspec_pullback_iso_construction}
  \lean{AlgebraicGeometry.Scheme.RelativeSpec.pullback_iso_construction}
  \uses{def:relspec_pullback_cocone, lem:relspec_pullback_iso_desc_isiso,
    lem:relspec_pullback_fst_isaffinehom}
  The canonical isomorphism of \(T\)-schemes
  \(T \times_X \underline{\Spec}_X(\mathcal{A}) \cong \underline{\Spec}_T(g^*\mathcal{A})\),
  defined as the inverse of the colimit descent map of the pullback cocone.
\end{definition}

\begin{proof}
  The descent map is an isomorphism by \cref{lem:relspec_pullback_iso_desc_isiso}, and
  the construction takes its inverse. Constructed directly in Lean.
\end{proof}

\begin{lemma}

[Existence of the base-change isomorphism]
  \label{lem:relspec_pullback_iso}
  \lean{AlgebraicGeometry.Scheme.RelativeSpec.pullback_iso}
  \uses{def:relspec_pullback_iso_construction}
  There exists a canonical isomorphism of \(T\)-schemes
  \(T \times_X \underline{\Spec}_X(\mathcal{A}) \cong \underline{\Spec}_T(g^*\mathcal{A})\).
  This is the existence statement consumed by the base-change theorem
  (\cref{thm:relative_spec_base_change}).
\end{lemma}

\begin{proof}

  Witnessed by the explicit construction
  \cref{def:relspec_pullback_iso_construction}. Proved directly in Lean.
\end{proof}

\section{Functoriality}
\label{sec:relspec_functoriality}

\begin{theorem}[Relative spectrum as a functor]
  \label{thm:relative_spec_functorial}
  \lean{AlgebraicGeometry.Scheme.RelativeSpec.functor}
  \uses{thm:relative_spec_univ, thm:relative_spec_affine_base}

  % NOTE (iter-173 review): the iter-173 file-skeleton encodes this as a
  % bare object-level function ``QcohAlgebra → Over X'' (not a categorical
  % ``Functor''), with morphism action + functoriality + the equivalence
  % onto affine X-schemes all deferred. Iter-174+ upgrades to the full
  % contravariant ``Functor (X.QcohAlgebra)ᵒᵖ (Over X)'' with the universal
  % property-induced ``map'' once QcohAlgebra is unpacked.
  % SOURCE: [Stacks Project], tag 01LR (definition-relative-spec +
  % lemma-glueing-gives-functor-spec + lemma-spec-properties part (3))
  % (read from references/stacks-constructions.tex, L602-L691).
  % SOURCE QUOTE:
  % "Let \(S\) be a scheme. Let \(\mathcal{A}\) be a quasi-coherent sheaf of
  %  \(\mathcal{O}_S\)-algebras. The {\it relative spectrum of \(\mathcal{A}\) over
  %  \(S\)}, or simply the {\it spectrum of \(\mathcal{A}\) over \(S\)} is the scheme
  %  constructed in Lemma REF which represents the
  %  functor \(F\) (REF), see
  %  Lemma REF.
  %  We denote it \(\pi : \underline{\Spec}_S(\mathcal{A}) \to S\).
  %  The ``universal family'' is a morphism of \(\mathcal{O}_S\)-algebras
  %  \(\)
  %  \mathcal{A}
  %  \longrightarrow
  %  \pi_*\mathcal{O}_{\underline{\Spec}_S(\mathcal{A})}
  %  \(\)
  %  [...]
  %  The universal map
  %  \(\)
  %  \mathcal{A}
  %  \longrightarrow
  %  \pi_*\mathcal{O}_{\underline{\Spec}_S(\mathcal{A})}
  %  \(\)
  %  is an isomorphism of \(\mathcal{O}_S\)-algebras."
  \textit{Source: [Stacks Project], tag 01LR (definition-relative-spec +
  lemma-glueing-gives-functor-spec + lemma-spec-properties);
  [Hartshorne], II~Ex.~5.17(b).}
  The construction \(\mathcal{A} \mapsto \underline{\Spec}_X(\mathcal{A})\) assembles into
  a contravariant functor
  \[
    \underline{\Spec}_X : \mathrm{QcohAlg}(X)^{op} \longrightarrow \mathrm{AffSch}/X
  \]
  from quasi-coherent \(\mathcal{O}_X\)-algebras to affine \(X\)-schemes: a morphism of
  \(\mathcal{O}_X\)-algebras \(\psi : \mathcal{A} \to \mathcal{B}\) induces, by the universal
  property (\cref{thm:relative_spec_univ}) applied to the pair
  $(\pi_\mathcal{B}, \psi \circ \pi_\mathcal{B}^*) :
  \pi_\mathcal{B}^*\mathcal{A} \to \pi_\mathcal{B}^*\mathcal{B} \to \mathcal{O}_{\Spec\mathcal{B}}$,
  an \(X\)-morphism
  \(\underline{\Spec}_X(\psi) : \underline{\Spec}_X(\mathcal{B}) \to \underline{\Spec}_X(\mathcal{A})\).
  Functoriality (identity, composition) is direct from the universal property. The
  pushforward \(\pi_* \mathcal{O}_{\underline{\Spec}_X(\mathcal{A})} \cong \mathcal{A}\)
  exhibits \(\underline{\Spec}_X\) as an equivalence between \(\mathrm{QcohAlg}(X)^{op}\) and
  the full subcategory of affine \(X\)-schemes, with inverse \(f \mapsto f_* \mathcal{O}_Y\)
  (Stacks lemma-spec-properties part~(3)).
\end{theorem}

\begin{proof}
  
  Functoriality follows directly from \cref{thm:relative_spec_univ}: for
  \(\psi : \mathcal{A} \to \mathcal{B}\), the pair
  \((\pi_\mathcal{B} : \underline{\Spec}_X(\mathcal{B}) \to X, \;
  \pi_\mathcal{B}^*(\psi))\) defines an \(X\)-morphism
  \(\underline{\Spec}_X(\mathcal{B}) \to \underline{\Spec}_X(\mathcal{A})\) via the
  representing universal pair. Identity and composition follow from naturality of the
  representing bijection in \(\mathcal{A}\). For the equivalence claim, \(\underline{\Spec}_X\)
  is fully faithful because the isomorphism
  \(\pi_* \mathcal{O}_{\underline{\Spec}_X(\mathcal{A})} \cong \mathcal{A}\)
  (locally identifying with \(\Gamma(\Spec(\mathcal{A}(U)),\mathcal{O}) = \mathcal{A}(U)\)
  via \cref{thm:relative_spec_affine_base}) reconstructs \(\mathcal{A}\) from
  \(\underline{\Spec}_X(\mathcal{A})\). Essential surjectivity onto affine \(X\)-schemes uses
  the canonical morphism \(X \to \underline{\Spec}_X(\pi_*\mathcal{O}_X)\) of Stacks
  lemma-canonical-morphism; on an affine \(X\)-scheme \(\pi : Y \to X\), this canonical
  morphism is an isomorphism (affine-locally it is the identity
  \(\Spec\Gamma(\pi^{-1}(U), \mathcal{O}_Y) \to \Spec\Gamma(\pi^{-1}(U), \mathcal{O}_Y)\)).
\end{proof}

\section{Lean encoding}
\label{sec:relspec_lean_encoding}

The Lean target file is the new \texttt{AlgebraicJacobian/Picard/RelativeSpec.lean}.
The construction proceeds by gluing absolute affine pieces: for each affine open
\(U \subseteq X\) with \(U = \Spec R\) and \(A := \mathcal{A}(U)\) regarded as an
\(R\)-algebra, the local piece is \(\Spec A\) with structure morphism \(\Spec A \to \Spec R = U\)
induced by the algebra unit \(R \to A\). Mathlib provides \texttt{IsAffineOpen.toScheme} for
the affine-open inclusion and the absolute \texttt{Scheme.Spec} functor for the local
pieces; the transition isomorphisms come from quasi-coherence of \(\mathcal{A}\) via
\texttt{QuasiCoherent.tensorIso} (or the explicit
\(\mathcal{A}(U) \otimes_R S \xrightarrow{\sim} \mathcal{A}(V)\) for \(V = \Spec S \subseteq U\)).
The gluing backbone is \texttt{AlgebraicGeometry.Scheme.GlueData} (general construction) or
the convenience wrapper \texttt{AlgebraicGeometry.AffineScheme.glueOpens}; the cocycle
condition reduces to Stacks \texttt{lemma-transitive-spec}. The universal property
(\cref{thm:relative_spec_univ}) is expressed as a
\texttt{CategoryTheory.Functor.RepresentableBy} witness, providing the natural bijection
$\Hom_X(T, \underline{\Spec}_X(\mathcal{A})) \cong
\Hom_{\mathcal{O}_X\text{-alg}}(\mathcal{A}, g_*\mathcal{O}_T)$ as a
\texttt{Yoneda}-style structure. The base-change isomorphism
(\cref{thm:relative_spec_base_change}) and functoriality
(\cref{thm:relative_spec_functorial}) are derived from the universal property by
Yoneda. No new core axiom is required; the entire chapter is a definitional/structural
construction on top of existing Mathlib affine-scheme infrastructure.

\section{Out of scope}
\label{sec:relspec_out_of_scope}

This chapter packages \emph{only} the relative-spectrum construction and its three core
properties. The following downstream constructions live in separate sibling chapters and
are explicitly out of scope here:
\begin{itemize}
  \item \textbf{Relative Picard functor \(\Pic^\sharp_{C/k}\)} (A.1.c) --- consumed-by, not
    defined-here; its blueprint home is the planned chapter
    \texttt{Picard\_RelPicFunctor.tex}.
  \item \textbf{Line-bundle pullback on a product \(C \times_k T\)} (A.1.b) ---
    \texttt{Picard\_LineBundlePullback.tex} (planned).
  \item \textbf{Connectedness, dimensionality, smoothness of \(\underline{\Spec}_X(\mathcal{A})\)}
    --- downstream properties that depend on assumptions on \(\mathcal{A}\)
    (locally finite type, flat, etc.), not unconditional consequences of the construction.
  \item \textbf{Identity-component \(\Pic^0\) and abelian-variety structure} (A.3) --- a
    separate Route~A sub-row, fed by but not constructed in this chapter.
\end{itemize}
